\documentclass[10pt,a4paper]{amsart}
\usepackage{amsmath,amssymb,amsthm,graphicx}
% Retained preamble for the Stacks Project chapter {CH} (shared project preamble).
% The chapter file pinned-original-source.tex only carries \input{{preamble}} on its first line,
% so no preamble text lives inside the chapter itself.  The Stacks Project uses one single shared
% preamble for every chapter: the snapshot stages it three times (sources/stacks-chow/preamble.tex,
% sources/stacks-etale-cohomology/preamble.tex and focused-seeds/AGP/sources/stacks-snapshot-20261011/
% preamble.tex) and those copies are identical in size.  The block below is copied mechanically,
% line for line, from sources/stacks-chow/preamble.tex lines 14-18 and 154-217: the Xy-pic setup and
% the native theorem declarations together with the mathematical macros.  Nothing was added, removed
% or redefined; the wrapper class and the amsmath/amssymb/amsthm/graphicx base packages are supplied
% by the assembler.
% For commutative diagrams we use Xy-pic
% Chapter: stacks-cohomology.tex -- Cohomology of Sheaves.
\usepackage[all]{xy}

% We use 2cell for 2-commutative diagrams.
\xyoption{2cell}

% Theorem environments.
%
\theoremstyle{plain}
\newtheorem{theorem}[subsection]{Theorem}
\newtheorem{proposition}[subsection]{Proposition}
\newtheorem{lemma}[subsection]{Lemma}

\theoremstyle{definition}
\newtheorem{definition}[subsection]{Definition}
\newtheorem{example}[subsection]{Example}
\newtheorem{exercise}[subsection]{Exercise}
\newtheorem{situation}[subsection]{Situation}

\theoremstyle{remark}
\newtheorem{remark}[subsection]{Remark}
\newtheorem{remarks}[subsection]{Remarks}

\numberwithin{equation}{subsection}

% Macros
%
\def\lim{\mathop{\mathrm{lim}}\nolimits}
\def\colim{\mathop{\mathrm{colim}}\nolimits}
\def\Spec{\mathop{\mathrm{Spec}}}
\def\Hom{\mathop{\mathrm{Hom}}\nolimits}
\def\Ext{\mathop{\mathrm{Ext}}\nolimits}
\def\SheafHom{\mathop{\mathcal{H}\!\mathit{om}}\nolimits}
\def\SheafExt{\mathop{\mathcal{E}\!\mathit{xt}}\nolimits}
\def\Sch{\mathit{Sch}}
\def\Mor{\mathop{\mathrm{Mor}}\nolimits}
\def\Ob{\mathop{\mathrm{Ob}}\nolimits}
\def\Sh{\mathop{\mathit{Sh}}\nolimits}
\def\NL{\mathop{N\!L}\nolimits}
\def\CH{\mathop{\mathrm{CH}}\nolimits}
\def\proetale{{pro\text{-}\acute{e}tale}}
\def\etale{{\acute{e}tale}}
\def\QCoh{\mathit{QCoh}}
\def\Ker{\mathop{\mathrm{Ker}}}
\def\Im{\mathop{\mathrm{Im}}}
\def\Coker{\mathop{\mathrm{Coker}}}
\def\Coim{\mathop{\mathrm{Coim}}}

% Boxtimes
%
\DeclareMathSymbol{\boxtimes}{\mathbin}{AMSa}{"02}

%
% Macros for moduli stacks/spaces
%
\def\QCohstack{\mathcal{QC}\!\mathit{oh}}
\def\Cohstack{\mathcal{C}\!\mathit{oh}}
\def\Spacesstack{\mathcal{S}\!\mathit{paces}}
\def\Quotfunctor{\mathrm{Quot}}
\def\Hilbfunctor{\mathrm{Hilb}}
\def\Curvesstack{\mathcal{C}\!\mathit{urves}}
\def\Polarizedstack{\mathcal{P}\!\mathit{olarized}}
\def\Complexesstack{\mathcal{C}\!\mathit{omplexes}}
% \Pic is the operator that assigns to X its picard group, usage \Pic(X)
% \Picardstack_{X/B} denotes the Picard stack of X over B
% \Picardfunctor_{X/B} denotes the Picard functor of X over B
\def\Pic{\mathop{\mathrm{Pic}}\nolimits}
\def\Picardstack{\mathcal{P}\!\mathit{ic}}
\def\Picardfunctor{\mathrm{Pic}}
\def\Deformationcategory{\mathcal{D}\!\mathit{ef}}

% --- appended 2026-10-12 (Chief): the project own definitions of the reference/slogan/history
% display environments, copied mechanically from the shared Stacks preamble lines 8-12
% (\usepackage{verbatim} plus the three \newenvironment lines, which the chapter preamble keeps
% verbatim from the project preamble). Nothing added, removed or redefined: these environments
% carry no mathematical content and the project itself typesets them as comments. Without these
% three lines a counted statement whose own text uses them fails with Environment slogan undefined.
\usepackage{verbatim}
\newenvironment{reference}{\comment}{\endcomment}
\newenvironment{slogan}{\comment}{\endcomment}
\newenvironment{history}{\comment}{\endcomment}
\usepackage{cleveref}
\crefname{originallemma}{original lemma}{original lemmas}
\Crefname{originallemma}{Original lemma}{Original lemmas}
\makeatletter
\newlabel{lemma-Rhom-strictly-perfect}{{08DM}{source}{}{sourcepointer-1}{}}
\newlabel{lemma-Rhom-strictly-perfect@cref}{{[originallemma][0][]08DM}{[1][0][]source}{}{}{}}
\makeatother
\crefname{originallemma}{original lemma}{original lemmas}
\Crefname{originallemma}{Original lemma}{Original lemmas}
\makeatletter
\newlabel{lemma-diagonal-better}{{0BYR}{source}{}{sourcepointer-2}{}}
\newlabel{lemma-diagonal-better@cref}{{[originallemma][0][]0BYR}{[1][0][]source}{}{}{}}
\makeatother
\crefname{originallemma}{original lemma}{original lemmas}
\Crefname{originallemma}{Original lemma}{Original lemmas}
\makeatletter
\newlabel{lemma-local-actual}{{08C9}{source}{}{sourcepointer-3}{}}
\newlabel{lemma-local-actual@cref}{{[originallemma][0][]08C9}{[1][0][]source}{}{}{}}
\makeatother
\crefname{originalreference}{original reference}{original references}
\Crefname{originalreference}{Original reference}{Original references}
\makeatletter
\newlabel{homology-section-truncations}{{0118}{source}{}{sourcepointer-4}{}}
\newlabel{homology-section-truncations@cref}{{[originalreference][0][]0118}{[1][0][]source}{}{}{}}
\makeatother
\title{The hom complex of a strictly perfect complex into a K-flat complex is again K-flat}
\author{The Stacks Project Authors}
\date{stacks-project:cohomology.tex:local-snapshot-20261011; GNU FDL 1.2 or later}
\begin{document}
\maketitle
\noindent\textit{Editorial note.} All delivered mathematical text is the author's own, copied verbatim from the pinned Stacks Project chapter file cohomology.tex; no statement or proof step was rewritten, reordered or removed.\par\medskip
\noindent\textit{Editorial note.} The counted span cites five cross-chapter labels, each recorded in assembly.external\_references with its canonical tag as a navigation pointer.\par\medskip
\noindent\textit{Editorial note.} The proof cites no bibliography item inside the counted span.\par\medskip
\noindent\textit{Allowed standard prerequisites (editorial).}\begin{itemize}
\item Computation of R-Hom out of a strictly perfect complex and the comparison map for tensor products. Stacks Project, Cohomology of Sheaves, Lemma lemma-Rhom-strictly-perfect (tag 0A8D) and Lemma lemma-diagonal-better (tag 0A8K). Supplies the description of the hom complex and the comparison map gamma used in the proof.
\item K-flat complexes and acyclicity of tensor products. Stacks Project, Cohomology of Sheaves, Lemma lemma-local-actual (tag 0A8L). Gives the acyclicity of the total tensor product with a K-flat complex.
\end{itemize}\medskip
\ifdefined\hypertarget\hypertarget{sourcepointer-1}{}\fi
\noindent\textit{Source reference.} Original lemma 08DM: sources/stacks-chow/tags.txt, source lines 10642--10642. Cross-chapter standard prerequisite cited by the author inside the counted span; it is used as a pointer and its own proof is not part of this item.\par\medskip
\ifdefined\hypertarget\hypertarget{sourcepointer-2}{}\fi
\noindent\textit{Source reference.} Original lemma 0BYR: sources/stacks-chow/tags.txt, source lines 15024--15024. Cross-chapter standard prerequisite cited by the author inside the counted span; it is used as a pointer and its own proof is not part of this item.\par\medskip
\ifdefined\hypertarget\hypertarget{sourcepointer-3}{}\fi
\noindent\textit{Source reference.} Original lemma 08C9: sources/stacks-chow/tags.txt, source lines 10594--10594. Cross-chapter standard prerequisite cited by the author inside the counted span; it is used as a pointer and its own proof is not part of this item.\par\medskip
\ifdefined\hypertarget\hypertarget{sourcepointer-4}{}\fi
\noindent\textit{Source reference.} Original reference 0118: sources/stacks-chow/tags.txt, source lines 1354--1354. Cross-chapter standard prerequisite cited by the author inside the counted span; it is used as a pointer and its own proof is not part of this item.\par\medskip
\begin{lemma}
\label{lemma-Rhom-strictly-perfect-K-flat}
In the situation of Lemma \ref{lemma-Rhom-strictly-perfect}
if $\mathcal{F}^\bullet$ is K-flat, then $\mathcal{H}^\bullet$ is K-flat.
\end{lemma}

\begin{proof}
Observe that $\mathcal{H}^\bullet$ is simply the hom complex
$\SheafHom^\bullet(\mathcal{E}^\bullet, \mathcal{F}^\bullet)$
since the boundedness of the strictly prefect complex
$\mathcal{E}^\bullet$ insures that the products in the definition
of the hom complex turn into direct sums.
Let $\mathcal{K}^\bullet$ be an acyclic complex of
$\mathcal{O}_X$-modules. Consider the map
$$
\gamma :
\text{Tot}(\mathcal{K}^\bullet \otimes
\SheafHom^\bullet(\mathcal{E}^\bullet, \mathcal{F}^\bullet))
\longrightarrow
\SheafHom^\bullet(\mathcal{E}^\bullet,
\text{Tot}(\mathcal{K}^\bullet \otimes \mathcal{F}^\bullet))
$$
of Lemma \ref{lemma-diagonal-better}. Since $\mathcal{F}^\bullet$ is K-flat,
the complex $\text{Tot}(\mathcal{K}^\bullet \otimes \mathcal{F}^\bullet)$
is acyclic, and hence by Lemma \ref{lemma-local-actual}
(or Lemma \ref{lemma-Rhom-strictly-perfect} if you like)
the target of $\gamma$ is acyclic too. Hence to prove the lemma it suffices to
show that $\gamma$ is an isomorphism of complexes. To see this, we may argue
by induction on the length of the complex $\mathcal{E}^\bullet$.
If the length is $\leq 1$ then the $\mathcal{E}^\bullet$
is a direct summand of $\mathcal{O}_X^{\oplus n}[k]$ for some
$n \geq 0$ and $k \in \mathbf{Z}$ and in this case the result
follows by inspection. If the length is $> 1$, then we reduce
to smaller length by considering the termwise split short exact
sequence of complexes
$$
0 \to \sigma_{\geq a + 1} \mathcal{E}^\bullet \to
\mathcal{E}^\bullet \to
\sigma_{\leq a} \mathcal{E}^\bullet \to 0
$$
for a suitable $a \in \mathbf{Z}$, see
Homology, Section \ref{homology-section-truncations}.
Then $\gamma$ fits into a morphism of termwise split
short exact sequences of complexes.
By induction $\gamma$ is an isomorphism for
$\sigma_{\geq a + 1} \mathcal{E}^\bullet$
and $\sigma_{\leq a} \mathcal{E}^\bullet$ and hence the result for
$\mathcal{E}^\bullet$ follows. Some details omitted.
\end{proof}

\end{document}
